% Part: incompleteness
% Chapter: introduction
% Section: decidability

\documentclass[../../../include/open-logic-section]{subfiles}

\begin{document}

\olfileid{inc}{int}{dec}

\olsection{عدم قابلية القرار وعدم الاكتمال}

يتطلب برهان غودل لمبرهنتي عدم الاكتمال حوسبة التركيب. لكن يمكننا،
حتى من دون ذلك، الحصول على نتائج حسنة بمجرد افتراض أن النظرية
!!{represents} جميع العلاقات !!{decidable}. والبرهان حجة قطرية تشبه
برهان عدم قابلية تقرير مسألة التوقف.

\begin{thm}
إذا كانت $\Gamma$ نظرية متسقة !!{represents} كل علاقة !!{decidable}،
فإن $\Gamma$ ليست !!{decidable}.
\end{thm}

\begin{proof}
افترض أن $\Gamma$ !!{decidable}. سنبين أنها، إذا كانت !!{represents}
كل علاقة !!{decidable}، لا بد أن تكون غير متسقة.

تمثل الخواص !!^{decidable} (أي العلاقات الأحادية) ب!!{formula}s ذات
متغير حر واحد. ولتكن $!A_0(x)$، و$!A_1(x)$، و\dots تعدادًا قابلًا
للحساب لجميع هذه ال!!{formula}s. انظر الآن إلى المجموعة الآتية
$D\subseteq\Nat$:
\[
D = \Setabs{n}{\Gamma \Proves \lnot !A_n(\num{n})}.
\]
المجموعة $D$ !!{decidable}؛ إذ يمكننا اختبار ما إذا كان $n\in D$
بحساب $!A_n(x)$ أولًا، ثم $\lnot !A_n(\num{n})$ منها. ومن الواضح أن
استبدال الحد $\num{n}$ بكل ورود حر لـ$x$ في $!A_n(x)$، ووضع
$\lnot$ قبل $!A_n(\num{n})$، أمر آلي. وبحسب الفرض تكون $\Gamma$
!!{decidable}، ولذلك يمكننا اختبار ما إذا كان
$\lnot !A_n(\num{n})\in\Gamma$. فإذا كان كذلك كان $n\in D$، وإلا كان
$n\notin D$. ومن ثم تكون $D$ أيضًا !!{decidable}.

ولما كانت $\Gamma$ !!{represents} جميع الخواص !!{decidable}، فهي
!!{represents}~$D$. وال!!{formula}s التي !!{represents}~$D$ في
$\Gamma$ كلها بين $!A_0(x)$، و$!A_1(x)$، و\dots. فليكن $d$ عددًا
بحيث !!{represents} $!A_d(x)$ المجموعة $D$ في~$\Gamma$. فإذا كان
$d\notin D$، لزم، لأن $!A_d(x)$ !!{represents}~$D$، أن
$\Gamma\Proves\lnot !A_d(\num{d})$. لكن هذا يعني أن $d$ يستوفي شرط
تعريف~$D$، ومن ثم $d\in D$، وهو يناقض $d\notin D$. ولذلك نستنتج
بطريق الخلف أن $d\in D$.

ولما كان $d\in D$، يفيد تعريف~$D$ أن
$\Gamma\Proves\lnot !A_d(\num{d})$. ومن جهة أخرى، لما كان
$!A_d(x)$ !!{represents}~$D$ في $\Gamma$، كان
$\Gamma\Proves !A_d(\num{d})$. ولذلك تكون $\Gamma$ غير متسقة.
\end{proof}

\begin{explain}
تبين المبرهنة السابقة أنه لا توجد نظرية متسقة !!{represents} جميع
العلاقات !!{decidable} وتكون مع ذلك !!{decidable}. وسنبين أن $\Th{Q}$
!!{represents} جميع العلاقات !!{decidable}؛ وهذا يعني أن جميع
النظريات التي تشتمل على $\Th{Q}$، مثل $\Th{PA}$ و$\Th{TA}$، تفعل
ذلك أيضًا، ومن ثم فهي أيضًا غير !!{decidable}. (ولما كانت هذه
النظريات كلها صادقة في النموذج القياسي، فهي كلها متسقة.)

ويمكننا أيضًا استعمال هذه النتيجة للحصول على صيغة ضعيفة من مبرهنة
عدم الاكتمال الأولى. فكل نظرية !!{axiomatizable} و!!{complete} تكون
!!{decidable}. ولذلك لا يمكن للنظريات المتسقة !!{axiomatizable} التي
!!{represents} جميع الخواص !!{decidable} أن تكون !!{complete}.
\end{explain}

\begin{thm}
إذا كانت $\Gamma$ !!{axiomatizable} و!!{complete}، فهي
!!{decidable}.
\end{thm}

\begin{proof}
كل نظرية غير متسقة !!{decidable}، لأن النظريات غير المتسقة تحتوي جميع
ال!!{sentence}s؛ ولذلك يكون جواب السؤال «هل $!A\in\Gamma$؟» دائمًا
«نعم»، أي يمكن تقريره.

افترض إذن أن $\Gamma$ متسقة، وأنها علاوة على ذلك !!{axiomatizable}
و!!{complete}. ولما كانت $\Gamma$ !!{axiomatizable} فهي
!!{computably enumerable}. إذ يمكننا تعداد جميع ال!!{derivation}s
الصحيحة من بديهيات~$\Gamma$ بدالة قابلة للحساب. ويمكننا أن نحسب من
!!{derivation} صحيح ال!!{sentence} التي !!{derive}ها؛ ولذلك توجد،
بالجمع بين الأمرين، دالة قابلة للحساب تعدد جميع مبرهنات~$\Gamma$.
وتكون !!{sentence} مبرهنة في~$\Gamma$ إذا وفقط إذا لم تكن
$\lnot !A$ مبرهنة، لأن $\Gamma$ متسقة و!!{complete}. ويمكننا إذن
تقرير ما إذا كانت $!A\in\Gamma$ كما يأتي. نعدد جميع مبرهنات
$\Gamma$. فإذا ظهرت $!A$ في القائمة عرفنا أن
$\Gamma\Proves !A$. وإذا ظهرت $\lnot !A$ فيها عرفنا أن
$\Gamma\Proves/ !A$. ولما كانت $\Gamma$ !!{complete}، فلا بد أن
تحصل إحدى الحالتين في نهاية المطاف، فينتج الإجراء جوابًا في النهاية.
\end{proof}

\begin{cor}
\ollabel{cor:incompleteness}
إذا كانت $\Gamma$ متسقة و!!{axiomatizable} و!!{represents} كل خاصية
!!{decidable}، فهي ليست !!{complete}.
\end{cor}

\begin{proof}
لو كانت $\Gamma$ !!{complete} لكانت !!{decidable} بالمبرهنة السابقة
(لأنها !!{axiomatizable} ومتسقة). لكنها، لما كانت !!{represents} كل
خاصية !!{decidable}، ليست !!{decidable} بالمبرهنة الأولى.
\end{proof}

\begin{prob}
بين أن $\Th{TA}=\Setabs{!A}{\Sat{N}{!A}}$ ليست
!!{axiomatizable}. ويجوز لك افتراض أن $\Th{TA}$ تمثل جميع الخواص
القابلة للقرار.
\end{prob}

حين نثبت أن $\Th{Q}$، مثلًا، !!{represents} جميع الخواص
!!{decidable}، تفيد النتيجة اللازمة أن $\Th{Q}$ لا بد أن تكون غير
تامة. لكن برهانها لا يقدم مثالًا ل!!{sentence} مستقلة؛ بل يبين مجرد
وجوب وجود !!a{sentence} كهذه. وللحصول على مثال لا بد أن نحوسب
التركيب ونتبع فكرة برهان غودل الأصلية. وبالطبع لا يزال علينا إثبات
الدعوى الأولى، وهي أن $\Th{Q}$ !!{represents} بالفعل جميع الخواص
!!{decidable}.

وينبغي التنبيه إلى أن كل نظرية \emph{مهمة} ليست بالضرورة غير تامة
أو غير قابلة للقرار. فثمة نظريات كثيرة قوية بما يكفي لوصف حقائق
رياضية مهمة، لكنها لا تستوفي شروط نتيجة غودل. فمثلًا
$\Th{Pres}=\Setabs{!A\in\Lang{L_{A^+}}}{\Sat{N}{!A}}$، أي مجموعة
ال!!{sentence}s في لغة الحساب من غير~$\times$ الصادقة في النموذج
القياسي، تامة وقابلة للقرار معًا. وتسمى هذه النظرية حساب بريسبرغر،
وهي تثبت جميع الحقائق عن الأعداد الطبيعية التي يمكن صوغها باستعمال
$\Obj{0}$ و$\prime$ و~$+$ وحدها.

\end{document}
